\documentclass[
    12pt,               % Tamanho da fonte
    a4paper,            % Tamanho do papel
    brazil              % Idioma principal do documento
]{abntex2}

% =================================================
% PACOTES ESSENCIAIS
% =================================================
\usepackage[utf8]{inputenc}
\usepackage[T1]{fontenc}
\usepackage{cmap}               % Mapeamento de fontes para busca no PDF
\usepackage{lmodern}            % Fonte Latin Modern
\usepackage{indentfirst}        % Indenta o primeiro parágrafo de cada seção
\usepackage{geometry}           % Controle de margens
\usepackage{microtype}          % Melhorias de alinhamento de texto

% Configuração de margens da ABNT
\geometry{
    top=3cm,
    bottom=2cm,
    left=3cm,
    right=2cm
}

% =================================================
% DADOS DO DOCUMENTO (Metadados para a Capa)
% =================================================
\instituicao{Instituto Federal de Alagoas -- IFAL \par Campus Maceió \par Bacharelado em Sistemas de Informação}
\autor{Emerson Gomes Vanderlei da Silva}
\titulo{Documento de Teste -- Ambiente LaTeX}
\local{Maceió}
\data{2026}
\preambulo{Documento de teste gerado para validação do ambiente de compilação local com suporte às normas da ABNT.}

% =================================================
% INÍCIO DO TEXTO
% =================================================
\begin{document}

% Gera a capa e a folha de rosto automaticamente
\imprimircapa
\imprimirfolhaderosto

% Sumário automático
\pdfbookmark[0]{\contentsname}{toc}
\tableofcontents*
\cleardoublepage

\section{Introdução}
Este é um documento de teste para verificar se o seu compilador LaTeX e a extensão \textit{LaTeX Workshop} no VS Code estão funcionando perfeitamente. 

O recuo do parágrafo acima está configurado de acordo com as normas da ABNT graças ao pacote \texttt{indentfirst}.

\section{Seção de Teste}
Você pode testar a acentuação e caracteres especiais, como ç, á, é, í, ó, ú, ã e õ. Se o PDF renderizar esses caracteres corretamente e abrir a visualização ao lado sem erros, o seu ambiente está 100\% pronto para produzir seus trabalhos acadêmicos!

\end{document}